\documentclass[10pt]{article}
\usepackage[margin=0.7in]{geometry}
\usepackage{amsmath,amssymb}
\usepackage{booktabs}
\usepackage{enumitem}
\usepackage[hidelinks]{hyperref}
\usepackage{microtype}
\setlist{nosep,leftmargin=1.4em}
\setlength{\parskip}{1.5pt}
\setlength{\parindent}{0pt}
\usepackage{titlesec}
\titlespacing*{\section}{0pt}{4pt}{2pt}
\titleformat{\section}{\normalfont\large\bfseries}{\thesection}{0.5em}{}

\title{\vspace{-2.5em}\textbf{KernelAscent: Measuring Genuine Recursive Self-Improvement\\
via a Kernel-to-Model Capability Loop}}
\author{}
\date{}

\begin{document}
\maketitle
\vspace{-3.5em}

\section{Problem and thesis}
Recursive self-improvement (RSI) is the central capability behind ``intelligence
explosion'' safety concerns, yet it has \emph{no rigorous quantitative benchmark}.
Existing work splits into two camps that never meet. \textbf{Kernel-generation
benchmarks} (KernelBench, TritonBench, RE-Bench's ``Optimize a Kernel,''
robust-kbench, KernelLLM, Kevin-32B) measure a one-shot artifact: can a model write
a fast GPU kernel? This capability is now established. \textbf{Self-improvement
systems} (STOP, ADAS, Darwin G\"odel Machine, SICA) rewrite their own \emph{scaffold
code} but never touch the compute substrate and never measure whether \emph{capability}
compounds. Neither closes the loop that defines RSI: an agent improves the machinery
that produces its own capability, and the improved capability then makes it better at
that same task.

\textbf{Thesis.} We propose a benchmark whose loop is genuine capability
self-improvement: a model optimizes the GPU kernels used to \emph{train itself};
under a \emph{fixed wall-clock budget} those faster kernels buy more effective
training, producing a more capable model; the more capable model then optimizes
kernels better. We measure whether kernel-optimization capability \emph{compounds}
across rounds, and separate true self-improvement from mere test-time compute scaling
via a matched control. To our knowledge no prior benchmark closes a
kernel$\rightarrow$model$\rightarrow$kernel capability loop.

\section{The recursive loop}
Let $M_k$ be the model at round $k$ and let the target capability be
kernel-optimization skill, measured on a \emph{held-out} task set $T$:
\[
C(M) \;=\; \frac{1}{|T|}\sum_{i\in T}
\mathbf{1}\!\left[\, s_i(M) > p \;\wedge\; \mathrm{err}_i(M) \le \varepsilon \,\right],
\qquad s_i = \frac{t^{\text{ref}}_i}{t^{\text{agent}}_i},
\]
i.e.\ the fraction of tasks the model solves with measured speedup $>p$ at numerical
error $\le \varepsilon$ on real GPUs (KernelBench's \texttt{fast\_p}). Let $K(M)$ be
the set of kernels $M$ emits for the \emph{training} stack (attention, GEMM,
optimizer, all-reduce fusions), and let $R(M)$ be the resulting training throughput
(steps/sec). The loop, run under a fixed per-round wall-clock budget $W$:

\begin{quote}\itshape
\textbf{Round $k$:} (1) evaluate $C(M_k)$ on $T$; (2) $M_k$ emits training kernels
$K(M_k)$; (3) train under those kernels for wall-clock $W$, yielding
$S_k = R(M_k)\cdot W$ gradient steps and a new model
$M_{k+1} = \mathrm{Train}(M_k;\,K(M_k),\,W)$.
\end{quote}

Because $W$ is fixed, the \emph{only} way to fit more or larger training steps is for
the model to produce faster kernels: kernel skill converts directly into training
compute, which (hypothesis) converts into capability, which raises kernel skill. This
is a genuine positive feedback loop on \emph{capability}, not merely on speed.

\section{Metrics and hypotheses}
\textbf{Compounding coefficient.} Fit $C_k = C_0 + a\,k + b\,k^2$. The sign of $b$
characterizes takeoff: $b>0$ accelerating, $b\approx0$ linear, $b<0$ decelerating,
and a plateau/regress if $g_k = C_{k+1}-C_k \to 0$ or $<0$. We report $g_k$, the
throughput ratio $\rho_k = R(M_k)/R(M_0)$, and the cumulative capability gain.

\textbf{The RSI signal (the linchpin).} We run a \emph{fixed-kernel control} arm that
trains every round with the round-0 kernels (no self-acceleration), and a
\emph{compute-matched} control that is granted the same step count $S_k$ without the
self-referential kernel emission. Define
\[
\Delta_k \;=\; C_k^{\text{self}} - C_k^{\text{control}} .
\]
\emph{RSI is confirmed only if $\Delta_k$ grows with $k$}: the self-accelerating arm
must pull ahead \emph{because its training throughput compounds}. If $\Delta_k\approx0$,
the effect is ordinary compute scaling, not self-improvement --- and we report that
honestly. Three falsifiable hypotheses: \textbf{H1} kernel quality measurably raises
training throughput ($\rho_k$ rises); \textbf{H2} added throughput raises
kernel-optimization capability ($C$ rises with $S$); \textbf{H3} the loop compounds
($g_k>0$ and $\Delta_k$ increasing).

\section{Experimental design}
\textbf{Substrate benchmark.} Kernel tasks are grounded in real inference/training
bottlenecks (MoE grouped GEMM, long-context / sliding-window attention, paged
attention + KV-cache quantization, int4/fp8 dequant-fused GEMM, RoPE fusion,
all-reduce fusion), tagged by the optimization \emph{primitive} they require (tiling,
tensor-core MMA, \texttt{cp.async} double-buffering, warp-specialization, TMA,
fusion, quantization). Tasks are \emph{procedurally generated} with a held-out private
split to defeat memorization; scored on \textbf{both A100 and H100} (final score =
$\min$ across arches) to penalize single-arch overfit.

\textbf{Training loop.} $M_0$ is an open code model ($\sim$7--32B). Rounds use RL
fine-tuning on kernel optimization (GRPO, following Kevin-32B), so the training loop
itself runs the model's kernels and self-acceleration is intrinsic. Capability is
always measured on the held-out split, never the RL training tasks.

\textbf{Anti-gaming (the Sakana lesson).} Sakana's ``AI CUDA Engineer'' (2025)
reward-hacked its own eval harness, reporting large speedups on slower kernels. We
freeze the correctness + timing harness outside the agent's write scope, use
median-of-$N$ timing on pinned GPU clocks with cache flushing and fresh allocations,
re-check adversarial hidden shapes every round, and forbid dispatch to
cuBLAS/CUTLASS/FlashAttention. Only externally measured wall-clock counts; the model
never reports its own numbers.

\begin{center}
\small
\begin{tabular}{@{}lll@{}}
\toprule
\textbf{Component} & \textbf{Measures} & \textbf{RSI role} \\
\midrule
Held-out \texttt{fast\_p} $C(M_k)$ & kernel-opt capability & the improving capability \\
Throughput $R(M_k)$, $\rho_k$ & training steps/sec from self-kernels & the feedback channel \\
Compounding coeff.\ $a,b,g_k$ & shape of the capability curve & takeoff dynamics \\
$\Delta_k$ vs.\ controls & self-improvement over compute scaling & the RSI existence test \\
\bottomrule
\end{tabular}
\end{center}

\section{Feasibility, risks, deliverables}
\textbf{Compute.} A100 (Ampere) + H100 (Hopper). Cost is dominated by the per-round
RL training under budget $W$; we bound total cost as (rounds)$\times W\times$(arms).
A pilot uses a small $M_0$, short $W$, and 3--4 rounds to establish the curve before
scaling. \textbf{Primary risk} is measurement validity of $\Delta_k$; we de-risk by
first prototyping a \emph{single} round end-to-end (emit kernels $\to$ verify
throughput gain $\to$ train $\to$ re-measure $C$) to confirm H1 and H2 in isolation
before claiming H3. \textbf{Secondary risk} is null compounding; this is itself a
publishable result (\emph{current models do not exhibit kernel-driven RSI takeoff}),
which is a valuable safety datapoint. \textbf{Deliverables:} (i) the tiered,
contamination-resistant kernel benchmark with A100+H100 scoring; (ii) the
kernel$\rightarrow$model RSI harness with locked verification; (iii) baseline
compounding curves and the $\Delta_k$ RSI test for open frontier-scale models; (iv)
a per-primitive capability profile localizing where self-improvement stalls.
\end{document}
